\chapter{Lateral Epicondylitis (Tennis Elbow)}
\index{Lateral Epicondylitis (Tennis Elbow)}
\label{sec:Lateral Epicondylitis (Tennis Elbow)}

\section{Overview}
Lateral epicondylitis, commonly known as tennis elbow, is a tendinopathy of the common extensor tendon origin at the lateral epicondyle of the humerus. It is an overuse injury characterized by pain and tenderness over the lateral elbow that is exacerbated by wrist extension activities. Despite the name, it frequently occurs in individuals who do not play tennis.

\section{Classification}
\begin{description}[font=\normalfont\bfseries,leftmargin=3em,labelwidth=2.5em]
  \item[Cause] Repetitive microtrauma from overuse of wrist extensor muscles
  \item[Organ System] Bones/Musculoskeletal
  \item[Pathophysiology] Repeated eccentric or concentric loading of the wrist extensors (primarily extensor carpi radialis brevis) causes collagen fiber disruption, tendinosis (angiofibroblastic hyperplasia), and failed healing response. The condition is primarily degenerative rather than inflammatory, despite the "-itis" suffix.
  \item[Duration] Acute to chronic; often persistent
  \item[Onset] Gradual
  \item[Mode of Transmission] Not applicable
  \item[Clinical Course] Slowly progressive lateral elbow pain that worsens with activity; may become chronic without intervention.
  \item[Severity] Mild to moderate; can be disabling in severe cases
  \item[Extent of Disease] Localized to the lateral elbow
  \item[Age Group] Most common in adults aged 35--55 years
  \item[Epidemiology] Affects 1--3\% of the general population annually. Peak incidence in the fifth decade. Equally common in men and women. Occupational risk is higher than sports-related risk.
  \item[ICD-11 Code] FA51.0
\end{description}

\section{Causes}
Lateral epicondylitis is caused by repetitive, forceful, or sustained use of the wrist extensor muscles, leading to microtears and degenerative changes at the common extensor tendon origin. Common culprits include occupational activities (computer work, carpentry, painting) and sports (tennis, squash, racquetball).

\section{Risk Factors}
\begin{itemize}[noitemsep]
  \item Occupational tasks involving repetitive wrist extension (keyboard use, assembly line work, plumbing)
  \item Racquet sports with poor technique or inappropriate equipment
  \item Age 35--55 years
  \item Smoking
  \item Obesity
  \item Prior upper extremity injury
\end{itemize}

\section{Signs and Symptoms}
\begin{itemize}[noitemsep]
  \item Pain and tenderness over the lateral epicondyle
  \item Pain aggravated by wrist extension against resistance
  \item Weak grip strength
  \item Pain when lifting objects or shaking hands
  \item Morning stiffness of the elbow
  \item Difficulty performing daily activities
\end{itemize}

\section{Progression and Stages}
The condition begins with mild pain after activity (reactive tendinopathy) and progresses to pain during activity that may become constant (tendinosis disrepair). In chronic, untreated cases, degenerative tendinopathy with structural tendon changes occurs. Symptoms typically persist for 6--24 months; up to 80\% resolve within one year with or without treatment.

\section{Complications}
\begin{itemize}[noitemsep]
  \item Chronic lateral elbow pain refractory to conservative therapy
  \item Reduced grip strength and functional impairment
  \item Tendon rupture (rare)
  \item Prolonged work disability or lost athletic participation
  \item Reduced quality of life
  \item Emotional and psychological effects
\end{itemize}

\section{Diagnosis}
\begin{itemize}[noitemsep]
  \item Clinical examination with palpation of lateral epicondyle
  \item Provocative tests (Cozen test, Mills test, chair lift test)
  \item Lab tests (not needed)
  \item Imaging tests (ultrasound or MRI for refractory cases to assess tendon integrity)
\end{itemize}

\section{Prevention}
\begin{itemize}[noitemsep]
  \item Use ergonomic equipment and proper technique for occupational and sports activities
  \item Strengthen forearm extensor and flexor muscles
  \item Warm up before and stretch after repetitive activities
  \item Regular check-ups
\end{itemize}

\section{Treatment Options}
\begin{itemize}[noitemsep]
  \item Activity modification and relative rest
  \item Physical therapy with eccentric wrist extensor exercises
  \item NSAIDs for pain relief
  \item Counterforce brace or epicondylar clasp
  \item Corticosteroid injection (short-term benefit)
  \item Extracorporeal shockwave therapy or platelet-rich plasma injection
  \item Surgery (open or arthroscopic release) for refractory cases
  \item Lifestyle changes
\end{itemize}

\section{Living with It}
\begin{itemize}[noitemsep]
  \item Follow treatment plan
  \item Regular appointments
  \item Adjust daily habits
  \item Support groups
  \item Keep learning
\end{itemize}

\section{Outlook}
Lateral epicondylitis is self-limiting in most patients, with 80--90\% recovering within 12--18 months regardless of treatment. Conservative therapy accelerates symptom improvement. Surgery is reserved for refractory cases and has a success rate of 80--90\%. Recurrence is possible, particularly if risk factors and precipitating activities are not addressed.
